\section{Instrumentos de medida distintos: el sesgo invisible}\label{sec:instrumentos}

Una variable ausente en un canal es el caso más visible de un problema general: cada canal mide con su propio instrumento. Con la notación de la sección~\ref{sec:canal}, sea $X^*$ la covariable latente, con valores en un espacio estándar de Borel $\X^*$, y sean $P^*$ y $Q^*$ sus leyes en el grupo objetivo y en el canal donante. Cada canal registra $X=\varphi_D(X^*)$, con $\varphi_0,\varphi_1\colon\X^*\to\X$ medibles, de modo que $P=P^*\circ\varphi_1^{-1}$ y $Q=Q^*\circ\varphi_0^{-1}$. Sea $\mu^*(x^*)=\E[Y(0)\mid X^*=x^*,D=0]$.

\begin{teorema}[Instrumentos de medida distintos]\label{teo:medicion}
Supongamos (A1), soporte latente $P^*\ll Q^*$ y (A2) latente: $\E[Y(0)\mid X^*,S=1]=\mu^*(X^*)$ $P^*$-casi seguro.
\begin{enumerate}
  \item \emph{Variables comunes.} Para toda aplicación medible $\xi$ sobre $\X^*$, $P^*\circ\xi^{-1}\ll Q^*\circ\xi^{-1}$. En particular, si una $\upsilon$ medible sobre $\X$ cumple $\upsilon\circ\varphi_0=\upsilon\circ\varphi_1$ (variables que ambos canales miden igual), las leyes de $W=\upsilon(X)$ en los dos grupos cumplen $P_W\ll Q_W$.
  \item \emph{Identificación con las variables comunes.} Si además $\mu^*(X^*)=\E[Y(0)\mid W,D=0]$ $Q^*$-casi seguro, los Teoremas~\ref{teo:identificacion}--\ref{teo:lebesgue} se aplican con $W$ como covariables.
  \item \emph{Medición distinta con soporte.} Si $\varphi_0$ es inyectiva y $P\ll Q$ (en espacios estándar de Borel, $\varphi_0$ tiene entonces imagen boreliana e inversa medible por el teorema de Lusin--Souslin \citep[teorema 15.1]{kechris1995}, y $\varphi_1(X^*)$ está en esa imagen $P^*$-casi seguro porque $Q$ se concentra en ella), el estimador ponderado con $X$ converge, con $Y$ acotado, a
  \[
  \theta_X=\int_{\X^*}\mu^*\bigl(\varphi_0^{-1}(\varphi_1(x^*))\bigr)\,P^*(dx^*),\qquad
  \theta_X-\theta=\int_{\X^*}\bigl[\mu^*\circ\varphi_0^{-1}\circ\varphi_1-\mu^*\bigr]\dd P^*.
  \]
  \item \emph{Indistinguibilidad.} En las condiciones de 3, existe un mundo sin error de medición, con $\varphi_1=\varphi_0$ y la ley latente del grupo objetivo igual a $P^*\circ(\varphi_0^{-1}\circ\varphi_1)^{-1}$, que reproduce la ley de los datos del grupo objetivo y del canal donante (y la de todo lo observado si, además, las covariables registradas de todas las unidades premium tienen una ley absolutamente continua respecto de $Q$), cumple los supuestos mantenidos y en el que $\theta=\theta_X$. Ningún procedimiento basado solo en los datos observados distingue $\theta$ de $\theta_X$.
\end{enumerate}
\end{teorema}
\begin{proof}
(1) Si $Q^*(\xi^{-1}(B))=0$, por el soporte latente $P^*(\xi^{-1}(B))=0$. Si $\upsilon\circ\varphi_0=\upsilon\circ\varphi_1$, entonces $W=\upsilon(\varphi_D(X^*))=(\upsilon\circ\varphi_0)(X^*)$ en los dos canales, y basta aplicar lo anterior con $\xi=\upsilon\circ\varphi_0$.

(2) Es la Proposición~\ref{prop:descartar} con $V=W$, que es función de $X^*$.

(3) En el canal donante, $X=\varphi_0(X^*)$ y $\varphi_0$ es inyectiva con inversa medible, así que condicionar en $X$ equivale a condicionar en $X^*$; por (A1), $m_0(x)=\mu^*(\varphi_0^{-1}(x))$ $Q$-casi seguro, y como $P\ll Q$ esta versión queda fijada $P$-casi seguro. El estimador ponderado converge a $\int m_0\dd P$ (Teorema~\ref{teo:identificacion}, que para esta convergencia solo usa $P\ll Q$), y por el cambio de variables de la medida imagen, $\int m_0\dd P=\int m_0(\varphi_1(x^*))\,P^*(dx^*)=\theta_X$. Por (A2) latente y la torre, $\theta=\int\mu^*\dd P^*$.

(4) En el mundo alternativo, el canal donante no cambia. El grupo objetivo se registra con $\varphi_0$ y su ley latente es $\tilde P^*=P^*\circ(\varphi_0^{-1}\circ\varphi_1)^{-1}$, así que su ley observada es $P^*\circ(\varphi_0\circ\varphi_0^{-1}\circ\varphi_1)^{-1}=P^*\circ\varphi_1^{-1}=P$, la misma que en el mundo real. Para su resultado observado basta tomar, dado $X$, la misma ley condicional que en el mundo real, y con (A2) latente su respuesta en el canal básico es $\mu^*$; su $\theta$ es $\int\mu^*\dd\tilde P^*=\theta_X$. El mundo alternativo cumple los supuestos: como $\varphi_0$ es inyectiva, $Q^*(A)=Q(\varphi_0(A))$ y $\tilde P^*(A)=P(\varphi_0(A))$, así que $P\ll Q$ da $\tilde P^*\ll Q^*$, y (A2) latente vale por construcción. Con la condición adicional, las unidades premium que no están en el grupo objetivo se tratan igual, transformando su ley latente por $\varphi_0^{-1}\circ\varphi_1$, que entonces está definida casi seguro; sin ella, $\varphi_1(X^*)$ puede caer fuera de la imagen de $\varphi_0$ para esas unidades, y un mundo con $\varphi_1=\varphi_0$ no puede reproducirlas.
\end{proof}

\begin{observacion}[Instrumentos aleatorios]\label{rem:stochastic}
Los puntos 3 y 4 usan el instrumento del grupo objetivo solo a través de la ley $P$ de sus covariables registradas. Por tanto, valen literalmente cuando ese instrumento es aleatorio, $X=\varphi_1(X^*,\xi)$ con $\xi$ una variable aleatoria auxiliar, como el error de una declaración: si $\varphi_0$ es inyectiva y $P\ll Q$, el estimador ponderado converge a $\theta_X=\E[\mu^*(\varphi_0^{-1}(X))\mid S=1]$, y el mundo en el que las covariables latentes del grupo objetivo se distribuyen como $\varphi_0^{-1}(X)$ dado $S=1$ reproduce los datos del grupo objetivo y del canal donante. No hace falta ninguna condición sobre cómo se relaciona $\xi$ con $Y(0)$, porque el límite depende del resultado solo a través de $m_0$, que fija el canal donante.
\end{observacion}

\paragraph{Lectura.} Una variable medida de forma distinta es más peligrosa que una ausente. La ausente rompe el soporte y normalmente deja huella: el clasificador ideal está seguro, $p=1$, y $\hat\varepsilon>0$ (la excepción, un tope en una de las fuentes, se comenta en la sección~\ref{sec:lalonde}). La medida de forma distinta puede conservar el soporte, y el clasificador interpreta la diferencia de medida como una diferencia real entre unidades. En un ejemplo con $P^*=N(1;\,1)$, $Q^*=N(0;\,1)$, $\mu^*(x)=\sigma(-3-0,8x)$, con $\sigma(t)=1/(1+e^{-t})$, y un canal premium que registra la variable desplazada en $\Delta$, un desplazamiento de $\Delta=\pm0,3$ desviaciones típicas produce sesgos de $-20,5\,\%$ y $+25,4\,\%$ sobre $\theta=2,92\,\%$. En una pequeña simulación de este ejemplo (detalles en el repositorio), el estimador ponderado con clasificador logístico coincide con $\theta_X$ ($z=+0,6$ en ambos casos). El AUC cambia (0,76 sin error; 0,82 y 0,69 con él), pero de una forma que también produciría una diferencia real de población. La única salida es información externa al diseño: unidades medidas con ambos instrumentos, por ejemplo clientes que cambiaron de canal, que permiten estimar $\varphi_0^{-1}\circ\varphi_1$, como en la calibración externa de \citet{webbvargas2017}.

\subsection{Evidencia con datos públicos: índice de masa corporal declarado frente a medido}\label{sec:nhanes}

El Teorema~\ref{teo:medicion} se puede comprobar directamente con datos en los que existen, para las mismas personas, los dos instrumentos.

\paragraph{Datos.} Usamos el ciclo 2017--2018 de la Encuesta Nacional de Salud y Nutrición de EE.\,UU. \citep{nchs2020}: los 4\,034 adultos de 20 años o más con talla y peso medidos y declarados, tensión arterial sistólica (media de hasta tres lecturas), perímetro de cintura, nivel educativo y ratio de renta sobre el umbral de pobreza. No usamos los pesos de la encuesta y tratamos la muestra como la población. El IMC declarado se calcula con el peso y la talla declarados y se redondea a 0,1, la precisión con la que se publica el IMC medido, para que sean comparables. El IMC declarado es de media 0,81 menor que el medido (desviación típica 1,99), y la diferencia crece con el IMC: de 0,00 en el cuartil más bajo del IMC medido a $-1,66$ en el más alto, en línea con la literatura sobre antropometría declarada \citep{connorgorber2007}. Usarlo solo en un canal es una medición distinta en el sentido del Teorema~\ref{teo:medicion}; que el error no sea clásico (crece con el IMC) es lo que hace que el sesgo sea apreciable.

\paragraph{Diseño.} Cada réplica asigna a cada persona al grupo objetivo (un canal premium sintético) con probabilidad $\sigma\bigl(-1+\beta_b z(\mathrm{IMC})+0,5\,z(\mathrm{edad})-0,4\,z(\mathrm{renta})\bigr)$, donde el IMC es el medido, $z$ estandariza y $\beta_b=0,8$; alrededor del 30\,\% acaba en el grupo objetivo. El canal es sintético, pero las covariables, los resultados y el error de medición son reales. Como la asignación depende solo de las covariables verdaderas, la intercambiabilidad latente y el soporte se cumplen por construcción, y $\theta$, el resultado medio del grupo objetivo, se observa en cada réplica. Comparamos cuatro especificaciones, todas con edad, sexo, raza, nivel educativo y renta: (A) IMC medido en los dos grupos; (B) IMC declarado en el grupo objetivo y medido en los donantes, el caso de medición por canal del Teorema~\ref{teo:medicion}; (C) sin IMC; (D) IMC declarado en los dos grupos. El estimador es el doblemente robusto con \emph{cross-fitting} en cinco pliegues, con un clasificador logístico penalizado y una regresión ridge, ambos aditivos en splines cúbicos; hay 200 réplicas. Las réplicas comparten la población y solo varía la asignación, así que evaluamos el sesgo, no la cobertura.

\paragraph{Predicciones registradas.} Antes de ejecutar calculamos el límite de la especificación B con el Teorema~\ref{teo:medicion}(3), con $\mu$, ajustada con toda la muestra, una estimación de $m_0(v,b)=\E[Y\mid V=v,\mathrm{IMC}=b,D=0]$, que aquí coincide con $\E[Y\mid V=v,\mathrm{IMC}=b]$ porque la asignación depende solo de $(V,\mathrm{IMC})$: promediando sobre las probabilidades de asignación $a_i$, el sesgo predicho es $\sum_ia_i[\mu(V_i,b_i^{\mathrm{decl}})-\mu(V_i,b_i)]/\sum_ia_i$, es decir, $-0,47$\,mmHg en la tensión y $-2,44$\,cm en la cintura. El valor declarado no es función del medido, así que este es el caso de instrumento aleatorio de la Observación~\ref{rem:stochastic}: como la asignación depende solo de $(V,\mathrm{IMC})$, se cumple la intercambiabilidad latente, $\theta=\int m_0(V,\mathrm{IMC})\dd P^*$, y el estimador converge a $\int m_0(V,\mathrm{IMC}^{\mathrm{decl}})\dd P^*$. El criterio registrado fue $|z|<3$, admitiendo un error de modelo de hasta el 25\,\% del sesgo predicho. Predijimos además ausencia de sesgo en A, ninguna masa de $\hat p$ cercana a 1 en A ni en B, una diferencia de AUC entre A y B inferior a 0,02 y, en C, un sesgo negativo y mayor en valor absoluto que el de B en la cintura.

\begin{table}[ht]
\centering
\caption{NHANES, diseño base (200 réplicas). Sesgo medio del estimador DR con su error de Monte Carlo entre paréntesis. En ninguna réplica hubo unidades objetivo con $\hat p>0,99$.}\label{tab:nhanes}
\begin{tabular}{@{}lcccc@{}}
\toprule
Especificación & Tensión (mmHg) & Cintura (cm) & AUC & $\neff/n_0$ \\
\midrule
A: medido en los dos & $-0,11$ (0,06) & $-0,18$ (0,02) & 0,743 & 0,41 \\
B: medición distinta & $-0,48$ (0,06) & $-2,56$ (0,02) & 0,715 & 0,50 \\
\quad predicción (Teorema~\ref{teo:medicion}) & $-0,47$ & $-2,44$ & & \\
C: sin IMC & $-1,78$ (0,05) & $-10,70$ (0,04) & 0,653 & 0,74 \\
D: declarado en los dos & $-0,47$ (0,06) & $-1,32$ (0,03) & 0,737 & 0,44 \\
\bottomrule
\end{tabular}
\end{table}

\paragraph{Resultados.} La tabla~\ref{tab:nhanes} confirma de forma aproximada la predicción del Teorema~\ref{teo:medicion}: $-0,48$ frente a $-0,47$\,mmHg ($z=-0,3$) y $-2,56$ frente a $-2,44$\,cm. La desviación en la cintura, del 5\,\%, está dentro de la tolerancia registrada, pero es estadísticamente distinta de cero ($z=-5,7$); descontando el pequeño sesgo de A, $-2,37$ frente a $-2,44$. Nada en los datos observados delata el problema. Ninguna unidad objetivo tiene $\hat p>0,99$, y $\neff/n_0$ \emph{mejora} de 0,41 a 0,50: el grupo objetivo, seleccionado por IMC alto, declara un IMC desplazado hacia abajo, hacia los donantes, así que los grupos se separan menos y los pesos parecen mejores justo cuando el estimador está sesgado (el ruido solo, en D, lleva $\neff/n_0$ únicamente a 0,44). El AUC baja de 0,743 a 0,715, pero la especificación A no está disponible en la práctica, así que esa bajada no se puede observar. Quitar el IMC es mucho peor ($-10,70$\,cm), porque el IMC determina la asignación. Fallan dos predicciones registradas: A tiene un sesgo de $-0,18$\,cm en la cintura ($z=-8,3$), que crece con la fuerza de la asignación y desaparece con LightGBM ($-0,01$; tabla~\ref{tab:nhanes-robust}), así que procede de los modelos de splines; y el AUC baja más que el umbral registrado de 0,02.

La especificación D, registrada como exploratoria, sugiere una lección práctica con límites: usar el IMC declarado en \emph{los dos} grupos reduce a la mitad el sesgo de la cintura respecto de B, pero no el de la tensión ($-0,47$ frente a $-0,48$). Medir los dos grupos con el mismo instrumento impreciso deja confusión residual, pero no evalúa el modelo de resultado en un argumento desplazado; cuál de los dos efectos domina depende del resultado. La tabla~\ref{tab:nhanes-robust} muestra que en la cintura D fue mejor que B en los cuatro diseños, entre un 24 y un 73\,\%, mientras que en la tensión fue peor en dos de ellos. No tenemos un resultado general sobre cuándo se da este orden.

\begin{table}[ht]
\centering
\caption{NHANES, robustez (200 réplicas cada uno). Sesgo medio del estimador DR en la cintura (cm) y, en las dos últimas columnas, en la tensión; entre paréntesis, la predicción registrada.}\label{tab:nhanes-robust}
\resizebox{\textwidth}{!}{%
\begin{tabular}{@{}lcccccc@{}}
\toprule
Diseño & A & B & C & D & B, tensión (mmHg) & D, tensión (mmHg) \\
\midrule
Asignación débil ($\beta_b=0,4$) & $-0,08$ & $-2,26$ ($-2,19$) & $-5,69$ & $-0,61$ & $-0,42$ ($-0,41$) & $-0,27$ \\
Base ($\beta_b=0,8$), splines & $-0,18$ & $-2,56$ ($-2,44$) & $-10,70$ & $-1,32$ & $-0,48$ ($-0,47$) & $-0,47$ \\
Asignación fuerte ($\beta_b=1,2$) & $-0,46$ & $-2,89$ ($-2,61$) & $-14,44$ & $-2,20$ & $-0,56$ ($-0,51$) & $-0,69$ \\
Base, LightGBM & $-0,01$ & $-2,42$ ($-2,44$) & $-10,38$ & $-1,23$ & $-0,20$ ($-0,47$) & $-0,45$ \\
\bottomrule
\end{tabular}}
\end{table}

\paragraph{Robustez.} Variamos la fuerza con la que el IMC determina la asignación y sustituimos los splines por LightGBM, registrando de nuevo las predicciones antes de ejecutar (tabla~\ref{tab:nhanes-robust}). Para cada diseño registramos cinco: B dentro de la tolerancia anterior; el sesgo de A en la cintura por debajo del 20\,\% del de B (un criterio relativo, elegido después de que la ejecución base mostrara un sesgo de modelo en A; con $|z|<3$, A también fallaría en los diseños débil y fuerte); ninguna unidad objetivo con $\hat p>0,99$ en A ni en B; C negativo y mayor que B en valor absoluto en la cintura; y D menor que B en valor absoluto en la cintura. De esas 15 predicciones se cumplen 14. En los cuatro diseños, ninguna unidad objetivo tiene $\hat p>0,99$ y $\neff/n_0$ es mayor en B que en A. El fallo es la tensión con LightGBM: $-0,20$ frente a $-0,47$ ($z=4,0$). La tensión depende poco del IMC y tiene mucho ruido, así que LightGBM aplana esa dependencia y solo transmite parte del error de medición, mientras que la predicción usó una $\mu$ ajustada con splines. Una comprobación no registrada lo respalda: con $\mu$ ajustada con LightGBM, el sesgo predicho es $-0,31$, y el $-0,20$ observado queda a $z=1,7$ (en la cintura, $-2,40$ en lugar de $-2,44$). Cuando el resultado depende poco de la covariable mal medida, la propia predicción depende del modelo de resultado.
